\documentclass{Thesis}
\usepackage[utf8]{inputenc}
\usepackage[T1]{fontenc}
\usepackage{geometry}
\usepackage{graphicx}
\usepackage{float}
\usepackage{hyperref}
\usepackage[english,indonesian]{babel}

\title{Optimasi Pemrosesan Template LaTeX Menggunakan Pendekatan Berbasis AI}
\author{Academic Support Agent}
\date{\today}

\begin{document}
\maketitle

\begin{abstract}

\end{abstract}

\chapter{PENDAHULUAN}


\chapter{TINJAUAN PUSTAKA}


% Original content (commented out):
% \documentclass[]{Thesis}
\usepackage[utf8]{inputenc}
\usepackage[english,mongolian]{babel}

\graphicspath{{./Pictures/}}
\begin{document}
\pagenumbering{roman}
\pagestyle{fancy}


\include{Information}

\input{Chapters/CoverPage}
\input{Chapters/TitlePage}
\input{Chapters/DefenceCom}
\input{Chapters/Abstract}

\tableofcontents

\addcontentsline{tof}{chapter}{List of Tables}
\listoffigandtable

\newpage
\chapter*{\listsymbolname}
\btypeout{\listsymbolname}
\addtotoc{\listsymbolname}
\input{Chapters/acrolist}


\pagenumbering{arabic}

\input{Chapters/Chapter1}
\input{Chapters/Chapter2}
\input{Chapters/Chapter3}
\input{Chapters/Chapter4} 
%\input{./Chapters/Chapter5} 
%\input{./Chapters/Chapter6} 
%\input{./Chapters/Chapter7} 

%---------
%	Ном зүй
%---------
\nocite{*} % Bibliography-д бичигдсэн бүх зүйлийг ном зүйд оруулна
\bibliographystyle{IEEEtran}
\bibliography{IEEEabrv,IEEEexample}
\appendix

\input{Chapters/appendix}

\end{document}


\end{document}